\section{Introduction}
\label{sec:introduction}

Neutral Higgs boson pairs of a two-Higgs-doublet model are produced at the LHC mainly by gluon fusion through top quark loops~\cite{Plehn:1996wb,Baglio:2023euv}. For pairs of pseudoscalars, $gg\to AA$, the next-to-leading order (NLO) QCD corrections were first computed in the limit of a heavy top quark~\cite{DawsonDittmaierSpira1998}, later matched to parton showers~\cite{Hespel:2014sla}, and recently obtained with the full dependence on the top quark mass~\cite{Baglio:2023euv}. In the benchmark point of Ref.~\cite{Baglio:2023euv} the full NLO cross section for $AA$ lies about $5\%$ below the Born-improved heavy top quark approximation at LHC energies, against about $12\%$ for $hH$, so the heavy top quark limit, with the exact mass dependence kept at leading order, is a reasonable starting point for higher orders. Beyond NLO only the two-loop virtual amplitudes in this limit are known~\cite{Bhattacharya:2019bmd}. For Standard Model Higgs pairs the corrections in the same limit are known to N$^3$LO~\cite{deFlorianMazzitelli2013NNLO,Banerjee:2018lfq,ChenEtAl2020}. Pairs of axion-like particles coupled to gluons are produced by the same mechanism, and non-resonant $gg\to aa$ production has recently been shown to constrain their interactions at the LHC~\cite{Brivio:2026alp}.

When the top quark is integrated out, the pseudoscalar couples to gluons through $G\tilde G$ and to light quarks through the divergence of the axial current. The amplitude for $gg\to AA$ then separates into two classes~\cite{Bhattacharya:2019bmd}. In Type~I both pseudoscalars originate from a single effective vertex. In Type~II each pseudoscalar is emitted from its own vertex and the two vertices are connected by gluon exchange in the $t$ or $u$ channel (Fig.~\ref{fig:diagrams}a). In the full theory the Type~II diagrams correspond to two top quark triangles joined by a gluon, the one-particle reducible contributions of Refs.~\cite{DawsonDittmaierSpira1998,Baglio:2023euv}; for Higgs pairs the corresponding class was computed at two loops in Ref.~\cite{Banerjee:2018lfq}. The two-loop Type~II finite remainders for $gg\to AA$ were obtained in closed form in Ref.~\cite{Bhattacharya:2019bmd}. They contain polylogarithms of weight up to four, among them $\Li_{2,2}$, and inverse powers of the Gram determinant $tu-\mA^4$, which vanishes at the pair threshold. The expressions occupy about \PublishedTwoLoopMB~MB, and large cancellations between their terms make them hard to evaluate close to threshold.

There are two reasons to study this region. The first is phenomenological. In the NLO predictions of Ref.~\cite{Baglio:2023euv} for $\mA=134.7$~GeV the invariant mass distribution of the pair is larger at $m_{AA}=300$~GeV, about 30~GeV above threshold, than at 400~GeV, and threshold resummation of this distribution requires at each $m_{AA}$ a hard function obtained from the infrared finite virtual amplitude~\cite{Sterman:1986aj,Catani:1989ne,deFlorian:2012za,deFlorian:2015moa,Bhattacharya:2021nsv}. The second concerns the analytic structure of loop amplitudes. When the transverse momentum $p_T$ of the final state particles vanishes, a one-loop amplitude with massless internal lines can have a Landau singularity~\cite{Landau:1959fi} in which each incoming parton splits into two collinear partons that then take part in two separate $2\to1$ scatterings. Gaunt and Stirling showed that this double parton scattering (DPS) singularity is at most logarithmic in $p_T$ for amplitudes of the Standard Model type, that for the processes they considered it requires the incoming partons to be in a state with $J_z=0$, and they gave a factorised formula for its coefficient~\cite{GauntStirling2011}. Dixon and Esterlis studied the corresponding singularity of six-gluon amplitudes in planar $\mathcal N=4$ super Yang--Mills theory to high loop order~\cite{DixonEsterlis2016}. As far as we are aware, it has not been examined beyond one loop in QCD at the level of amplitudes for a $2\to2$ process with a colour singlet final state; next-to-leading order corrections to the underlying splitting mechanism are known for double parton distributions~\cite{Diehl:2019rdh}. The Type~II amplitude has this topology: each incoming gluon can split into two gluons, which fuse pairwise into the two pseudoscalars at the effective vertices (Fig.~\ref{fig:diagrams}b). At threshold $p_T\to0$ at every scattering angle, so an expansion about threshold probes this region directly. In hadronic collisions the region in which both incoming partons split overlaps with double parton scattering in which both parton pairs arise from perturbative splittings, and the two contributions have to be combined without double counting~\cite{GauntStirling2011,DiehlGauntSchoenwald2017}; the related contribution with a splitting in one proton has been studied for pairs of Higgs bosons~\cite{GauntMaciulaSzczurek2014}.

In this paper we expand the two-loop Type~II finite remainders of Ref.~\cite{Bhattacharya:2019bmd} in the velocity $\bA$ of the pseudoscalars at fixed scattering angle, through $\bA^6$. Our main results are the following.
\begin{enumerate}
 \item After the endpoint constants of the expansion are expressed in a basis of \BasisCountWord{} constants, all inverse powers of $\bA$ and the odd powers through $\bA^5$ cancel, and every coefficient is even in $\cos\theta$, as required by Bose symmetry. At leading power the angle enters only through the logarithm of the transverse momentum of the pseudoscalars.
 \item This logarithm is the DPS singularity. At one loop its coefficient follows the factorised formula of Ref.~\cite{GauntStirling2011}, with the helicity weights of the splitting $g\to gg$ and the opposite signs of the amplitudes for $g^+g^+\to A$ and $g^-g^-\to A$. In the forward limit, where the momentum fractions depend on the velocity, and with the overall constant fixed at threshold, the formula reproduces the velocity dependence of the exact one-loop coefficient to \DPSMaxDeviation{} for $\DPSBetaMin\leq\bA\leq\DPSBetaMax$.
 \item At two loops the equal helicity amplitude grows as $\ln^3p_T$. The opposite helicity amplitude vanishes as $\bA^2$ at tree level and has no DPS logarithm at one loop, but grows as $\ln^2p_T$ at two loops. The ratio of the leading logarithms at two and one loops is the one found to all orders in planar $\mathcal N=4$ super Yang--Mills theory~\cite{DixonEsterlis2016}.
 \item In the finite virtual corrections to the cross section the one-loop logarithm cancels. At two loops $\ln^3p_T$ cancels, and the highest power is $\ln^2p_T$, with coefficient $-\tfrac74\pi^2\Nc^2$.
\end{enumerate}
We also find that the two-loop remainders released with Ref.~\cite{Bhattacharya:2019bmd} are larger by a factor of \FileFactor{} than the tree and one-loop amplitudes of the same reference, and that the released one-loop remainders omit their dependence on the renormalisation scale, so that they are correct only for $\muR=1$~GeV. We give all results in a consistent normalisation. The series reproduce the unexpanded remainders with relative errors of at most \PtwoErrLePointTwo{} for $\bA\leq\RecommendedBeta$; they are released together with an evaluator~\cite{SupplementalMaterial}.

The paper is organised as follows. Section~\ref{sec:setup} collects the conventions, the normalisation of the remainders and the counting of orders. Section~\ref{sec:expansion} describes the expansion and its accuracy. Section~\ref{sec:physics} discusses the leading power terms and their relation to the DPS singularity, and Sec.~\ref{sec:rate} their consequences for the cross section. We conclude in Sec.~\ref{sec:conclusions}.
